\documentclass{article}
\usepackage[left=2cm, right=2cm, top=2cm, bottom=2cm]{geometry}
\usepackage{graphicx}
\usepackage{amsmath}
\usepackage{amssymb}
\usepackage{amsthm}
\usepackage{fancyhdr}
\usepackage{verbatim}
\usepackage{listings}
\usepackage{xcolor}
\usepackage{pdfpages}
\usepackage{cancel}
\usepackage{physics}

\lstset{
    language=C++,
    basicstyle=\ttfamily\footnotesize,
    keywordstyle=\color{blue},
    commentstyle=\color{green},
    stringstyle=\color{red},
    numbers=left,
    numberstyle=\tiny\color{gray},
    frame=single,
    breaklines=true,
    captionpos=b,
}


\title{MTH 553: Lecture \# 21}
\author{Cliff Sun}

\newtheorem{theorem}{Theorem}[section]
\newtheorem{lemma}[theorem]{Lemma}
\newtheorem{definition}[theorem]{Definition}
\newtheorem{conjecture}[theorem]{Conjecture}
\newtheorem{proposition}[theorem]{Proposition}
\newtheorem{corollary}[theorem]{Corollary}
\newtheorem{one minute paper}[theorem]{One Minute Paper}

\pagestyle{fancy}
\lhead{\textbf{\thepage}\ \ \nouppercase{\rightmark}}
\chead{MTH 553: Lecture \# 21}
\rhead{Cliff Sun}

\begin{document}

\maketitle

\section*{Lecture Span}
\begin{enumerate}
    \item Extension from last time, heat kernel stuff
\end{enumerate}

\section*{Proof of some parts of Theorem from last time}

\begin{enumerate}
    \item True when $g \equiv 1$, then $u \equiv 1$. 
    \begin{proof}
        \begin{align*}
            \int_{\mathbb{R}^n}K(x,y,t)dy &= \int_{\mathbb{R}^n}\frac{1}{(4\pi t)^{n/2}}e^{-|x-y|^2/4t}dy \\
            &= \prod_{j=1}^{n}\frac{1}{(4\pi t)^{/2}}\int_{-\infty}^{\infty}e^{-(x_j - y_j)^2/4t}dy_j \\
            &= \prod_{j=1}^{n}\frac{1}{\sqrt{\pi}}\int_{-\infty}^{\infty}e^{-s_j^2}ds_j \\
            &= 1
        \end{align*}
    \end{proof}
    \item $u$ is bounded since $|u(x,t)| \leq ||g||_{\infty}$
    \item $u \in C^{\infty}(\mathbb{R}^n \times (0, \infty))$ by differentiating through the integral with respect to $x$ and $t$. This is true regardless if $g$ is smooth or not. 
    \item The heat kernel satisfies the heat equation. That is $K_t = \triangle K$. Therefore, $u_t = \triangle u$ by differentiating through the integral.
    \item Check initial condition, show that $u(x,t) \rightarrow g(\tilde{x})$ as $x \rightarrow \tilde{x}$ and $t \rightarrow 0$. 
    \item Uniqueness is guaranteed in bounded solution. 
\end{enumerate}

\subsection*{Remarks}

\begin{enumerate}
    \item $K(x,y,t)$ is the temperature of $x$ at time $t$ due to unit of heat released at $y$ at time $0$. 
    \item If $g(y) = \delta(y - \tilde{y})$, then graphically, you can see the heat spreading out from $\tilde{y}$. 
\end{enumerate}

\section*{Non-homogenous heat equation}

\begin{align*}
    u_t &= \triangle u + f(x,t)\\
    u(x,0) &= 0
\end{align*}

Duhamel also worked for the heat equation. Let $U(x,t;s)$ solve 

\begin{align*}
    U_t &= \triangle U \\
    U(x,0;s) &= f(x,s)
\end{align*}

Then, solve for $U$ using the heat kernel. Then, 

\begin{equation}
    u(x,t) = \int_{0}^{t}U(x,t-s;s) ds
\end{equation}

Solves the non-homogenous wave equation:

\begin{equation}
    u(x,t) = \int_{0}^{t}\int_{\mathbb{R}^n}K(x,y,t-s)f(y,s)dyds
\end{equation}

\begin{proof}
    \begin{align*}
        u_t &= \int_{0}^{t}U_t(t - s; s) ds + U(x,t-t;t) \\
        &= \int_{0}^{t} \triangle U(t-s;s)ds + f(x,t)\\
        &= \triangle u + f(x,t)
    \end{align*}
    The conversion of $U(x,0;t)$ to $f(x,t)$ is due to the initial conditions.
\end{proof}

\section*{Mean value formula for heat equation}

\begin{definition}
    For $x \in \mathbb{R}^n$, $t \in \mathbb{R}$, and $r > 0$. Define the heat ball as the following:
    \begin{equation}
        E(x,t;r) = \{(y,s) \in \mathbb{R}^n \times \mathbb{R}: s \leq t \text{ and } K(x,y,t-s) \leq 1/r^n\}
    \end{equation}
    This is not really a ball, but a closed elliptical curve. Note that the top point of the ball is the average of the heat inside the ball.
\end{definition}

\subsection*{Shape}

Take $x=0$, $t=0$. Let $E(r) = E(0,0;r)$ and $(y,s) \in E(r)$ such that $s \leq 0$ and $K(-y, -s) \leq 1/r^n$. Then we obtain 

\begin{equation}
    \frac{1}{(4\pi (-s))^{n/2}}e^{-|y|^2/4(-s)} \leq \frac{1}{r^n}
\end{equation}

\begin{equation}
    s \leq 0, \quad -\frac{n}{2}\log(-4\pi s) + \frac{|y|^2}{4s} + n\log r \geq 0
\end{equation}

Let 

\begin{equation}
    \Psi(y,s) = -\frac{n}{2}\log(-4\pi s) + \frac{|y|^2}{4s} + n\log r
\end{equation}

Continuing, 

\begin{equation}
    s \leq 0, \quad |y|^2 \leq 2n s \log(-4\pi s) - (4n\log(r))s
\end{equation}

This function is concave with respect to $s$. This is zero at $s=0$ and $s= -r^2/4\pi$. Then, the shape of the heat ball is 

\begin{equation}
    -\frac{r^2}{4\pi} \leq s \leq 0 
\end{equation}

And $$|y| \leq \sqrt{2ns\log(-4\pi s) - (4n\log (r))s}$$

The above is the cross section of the heat ball. 

\end{document}